% ==========================================================================
\chapter{绪论}
\pagestyle{fancy}
% ==========================================================================

\section{研究背景与意义}

常微分方程初值问题
\begin{equation}
  \frac{\dd y}{\dd t}=f(t,y),\qquad y(t_0)=y_0,
  \label{eq:ode-ivp}
\end{equation}
是数值分析中最基本的模型问题之一。无论是化学反应动力学、电路仿真、生物网络，还是偏微分方程经过空间离散后得到的半离散系统，最终都归结为形如 \eqref{eq:ode-ivp} 的初值问题。对于这类问题，数值方法的研究围绕三个相互牵制的目标展开：精度、稳定性和计算效率。显式方法（如显式 Euler、显式 Runge--Kutta）每步只需计算若干次右端函数，实现简单，但其稳定域有界；隐式方法（如隐式 Euler、Gear 方法）具有更好的稳定性，但每步需要求解线性或非线性方程组。当问题同时含有快、慢两种时间尺度时，这种矛盾尤为突出。

这种多尺度特征通常用“刚性”（stiffness）来刻画。粗略地说，若微分方程的解由缓慢变化的模态主导，而扰动中却含有衰减极快的模态，则问题是刚性的。以线性常系数系统 $y'=Ay$ 为例，若 Jacobian 矩阵 $A$ 的特征值实部均为负，且最大与最小模长之比很大，则称问题是刚性的。刚性带来的困难并不在于解的性态复杂，而在于显式方法的稳定步长被最快的衰减模态所限制，尽管该模态在解中几乎不可见。隐式方法虽然可以放宽这一限制，却牺牲了每步的计算简便性。

IMEX（implicit--explicit，隐式—显式）方法正是在这一背景下发展起来的一类分裂时间推进格式。其基本思想是把右端算子分裂为刚性部分与非刚性部分：对刚性部分采用隐式处理以保证稳定性，对非刚性部分采用显式处理以控制计算量。以本文考虑的半离散系统
\begin{equation}
  \frac{\dd \bU}{\dd t}=\bC(\bU)+\nu\bD\bU
  \label{eq:semidiscrete-intro}
\end{equation}
为例，其中 $\bC$ 为对流算子，$\nu\bD$ 为耗散算子，IMEX-Euler 格式为
\begin{equation}
  \frac{\bU^{n+1}-\bU^{n}}{\Delta t}
  =\bC(\bU^{n})+\nu\bD\bU^{n+1}.
  \label{eq:imex-euler-intro}
\end{equation}
每一步只需解一个线性方程组，既避免了对流项的隐式非线性迭代，又消除了扩散项的显式稳定性限制。这类格式在反应扩散方程、大气模拟、等离子体输运和 Kinetic 方程中都有广泛应用。

本文的研究对象是上述思想的一个具体而富有代表性的实例：一维粘性 Burgers 方程
\begin{equation}
  u_t+\left(\frac{u^2}{2}\right)_x=\nu u_{xx},\qquad \nu>0.
  \label{eq:burgers-intro}
\end{equation}
Burgers 方程是 Navier--Stokes 方程的一维简化模型，既保留了对流项的非线性与激波形成机制，又包含了扩散项的光滑化作用。它的粘性激波具有显式解析解，便于对数值格式进行定量检验。因此，它常被用作检验双曲型数值方法的“标准试金石”。

从常微分方程初值问题的角度看，对 \eqref{eq:burgers-intro} 作空间有限体积离散并采用直线法（method of lines, MOL），可得形如 \eqref{eq:semidiscrete-intro} 的大型刚性初值问题。其 Jacobian 的特征值中，对流部分量级为 $O(1/\Delta x)$，扩散部分量级为 $O(\nu/\Delta x^2)$。当网格加密时，扩散特征值的模迅速增大，刚性比随之发散。因此，Burgers 方程恰好提供了一个可以精确控制刚性强弱、又可以与解析解比较的试验平台。

本文不局限于某一个固定的分裂方式，而是引入分裂参数 $\theta\in[0,1]$，把耗散算子按比例分配到显式与隐式部分：
\begin{equation}
  \frac{\dd \bU}{\dd t}
  =\bC(\bU)+\theta\nu\bD\bU+(1-\theta)\nu\bD\bU,
  \label{eq:theta-family-intro}
\end{equation}
其中前两项（对流项与比例为 $\theta$ 的耗散项）作显式处理，最后一项（比例为 $1-\theta$ 的耗散项）作隐式处理。$\theta=1$ 对应耗散全显式，$\theta=0$ 对应耗散全隐式，$\theta=1/2$ 对应耗散对半分配。这个单参数族把“全显式”和“全隐式”两个端点用一条连续路径连接起来，使得我们可以系统地回答如下问题：在什么条件下，一个分裂是“容许”的？稳定性边界由什么决定？稳定性与精度之间如何权衡？

研究这一问题的意义有两方面。理论上，它把常微分方程稳定性理论中的 A-稳定、L-稳定等概念，推广到“双变量稳定性函数”的框架中，为分析分裂型 IMEX 格式提供了一个统一的几何图像。实践上，所得判据可以直接指导参数选择与时间步长控制，并为后续研究高维双曲守恒律、松弛系统和动理学方程提供可迁移的方法与代码基础。

\section{国内外研究现状}

\subsection{IMEX 时间推进方法}

IMEX 方法的思想可以追溯到 Crouzeix 在 1980 年提出的隐式—显式线性多步法\cite{crouzeix1980}。20 世纪 90 年代以后，随着对流扩散、反应扩散等问题的大规模计算需求，IMEX Runge--Kutta 方法得到系统发展。Ascher、Ruuth 和 Spiteri 在 1997 年构造了面向时变偏微分方程的一类 IMEX Runge--Kutta 格式，并给出阶条件与稳定性分析\cite{ascher1997}。Pareschi 和 Russo 将 IMEX Runge--Kutta 方法应用于带刚性松弛的双曲系统，强调了渐近保持性质\cite{pareschi2005}。Boscarino 则从微分代数方程的角度分析了 IMEX 方法的误差传播与阶退化现象\cite{boscarino2007}。2024 年，Boscarino 与 Russo 对拟线性双曲系统的渐近保持方法作了系统综述\cite{boscarino2024}，其中对源项分裂、刚性极限和容许性条件的讨论，构成了本文重要的方法论背景。

国内学者在 IMEX 方法及其应用方面也有大量工作，研究范围覆盖对流扩散方程、浅水方程、多相流和动理学方程等。总体而言，现有研究多针对某一具体格式或某一类方程，对“分裂参数如何影响稳定性”这一设计层面的问题，系统性的参数化分析仍不多见。

\subsection{刚性与稳定性理论}

刚性常微分方程的稳定性理论奠基于 Dahlquist 关于线性多步法 A-稳定的经典工作\cite{dahlquist1963}。A-稳定要求方法的绝对稳定域包含整个左半复平面；L-稳定进一步要求放大因子在 $\mathrm{Re}\,z\to-\infty$ 时趋于零，从而对最快的衰减模态具有“刚性衰减”作用。Hairer 与 Wanner 的专著\cite{hairer1996}对刚性问题的理论、算法与实现作了全面总结。Gear 的早期工作\cite{gear1971}则奠定了刚性积分器的实践基础。

对于分裂型 IMEX 格式，稳定性不再由单个复变量 $z$ 描述，而需要同时考虑显式部分与隐式部分的贡献。设 $\lambda_E$ 与 $\lambda_I$ 分别为显式、隐式算子的特征值，令 $z_E=\Delta t\lambda_E$，$z_I=\Delta t\lambda_I$，则 IMEX-Euler 的放大因子为
\begin{equation}
  R(z_E,z_I)=\frac{1+z_E}{1-z_I}.
  \label{eq:R-intro}
\end{equation}
这一双变量稳定性函数是本文的核心工具。Zakerzadeh 与 Noelle 在 2018 年针对刚性双曲守恒律的 IMEX 通量分裂，比较了对称分裂与非对称分裂的稳定性：对称分裂可以做到对刚性参数一致稳定，而非对称分裂必须逐个检验，他们并给出了存在小范围不稳定区域的反例\cite{zakerzadeh2018}。本文将这一思想具体化为耗散算子的按比例分裂，并给出精确的有限刚性稳定性边界。

\subsection{有限体积方法与直线法}

本文的空间离散采用一阶有限体积方法。有限体积方法以守恒形式为出发点，直接对单元平均值进行演化，天然保证离散守恒性。Godunov 在 1959 年提出的迎风格式\cite{godunov1959}和 Rusanov 在 1961 年提出的局部 Lax--Friedrichs 通量\cite{rusanov1961}，是两类最基础的数值通量。van Leer 的 MUSCL 重构\cite{vanleer1979}通过斜率限制器把格式推广到二阶，同时保持总变差有界（TVD）。LeVeque 的专著\cite{leveque2002}与 Toro 的专著\cite{toro2009}对有限体积方法的理论、通量构造和激波捕捉作了系统阐述。

直线法把空间离散与时间推进分离：先对半离散系统 \eqref{eq:semidiscrete-intro} 选用 ODE 求解器，再单独分析时间格式的稳定性。这种模块化结构正是本文代码设计的基础，也使得常微分方程初值问题的算法分析可以直接服务于偏微分方程的数值模拟。

\subsection{松弛系统与渐近保持格式}

Jin 与 Xin 在 1995 年提出了用于双曲守恒律的松弛格式\cite{jinxin1995}：用一个线性双曲系统加刚性源项逼近非线性守恒律，当松弛参数 $\varepsilon\to0$ 时恢复到原方程。该框架的优点是通量线性化，便于构造高分辨率格式；代价是源项具有 $1/\varepsilon$ 的刚性，必须采用隐式或渐近保持（asymptotic preserving, AP）处理。Jin 随后系统发展了 AP 格式的理论\cite{jin1999}，Dimarco 与 Pareschi 则对动理学方程的数值方法作了综述\cite{dimarco2014}。

在松弛系统中，源项分裂方式直接决定了格式的稳定性与渐近相容性。全隐式源项通常具有最好的稳定性，但每步需要求解耦合方程；只隐式化线性刚性部分、把非线性部分显式处理，则可以降低实现难度，却可能破坏容许性。本文将在第六章用双变量稳定性函数的语言，对这两种分裂作系统比较。

\section{本文主要工作}

围绕“参数化分裂族的耦合刚性稳定性”这一主题，本文完成了以下工作。

第一，建立了一个覆盖全显式到全隐式的单参数分裂族。引入 $\theta\in[0,1]$，把耗散算子按 $\theta$ 与 $1-\theta$ 分配至显式与隐式部分，并以 IMEX-Euler 作为时间推进格式。该族把常见的三种分裂（全显式、对半、全隐式）纳入同一分析框架。

第二，定义了双变量稳定性函数 $R(z_E,z_I)$，并给出了耦合刚性极限的精确刻画。对线性化模型 $u_t+au_x=\nu u_{xx}$，在耗散主导、对流可忽略的极限下，证明 $|R|\le1$ 当且仅当 $\theta\le1/2$；更精确地，当 $\theta>1/2$ 时，稳定当且仅当 $|z_D|\le2/(2\theta-1)$。本文还给出两种证明：代数证明与复平面圆盘的几何证明。进一步区分了 A-稳定与 L-稳定：只有 $\theta=0$ 具有刚性衰减，$\theta=1/2$ 是临界情形。

第三，完成了离散 Fourier 符号分析。对一阶 Rusanov 通量，导出放大因子并给出 CFL 型条件；通过网格尺度（$\xi=\pi$）分析，解释了 $\theta>1/2$ 时最大稳定步长按 $O(\Delta x^2/\nu)$ 收缩的机制。分析同时表明，Rusanov 通量自带的数值扩散会略微放宽稳定域，但在 $\Delta x\to0$ 时该效应消失，理论边界是尖锐的。

第四，设计并完成了四组数值实验：固定 CFL 下的刚性相图、不同 $\theta$ 下的最大稳定步长扫描、粘性激波的精度测试，以及网格收敛性测试。实验采用一阶 Rusanov 通量与 MUSCL/minmod 重构两种空间格式，输出激波位置误差、过渡层厚度和 $L^1$ 误差。实验结果与理论预测一致，并揭示了“稳定”与“准确”并不等价：$\theta=1/2$ 是稳定性的临界点，却不是精度最优的选择。

第五，将上述框架扩展到 Burgers 方程对应的 Jin--Xin 松弛系统。比较了源项全隐式（策略 A）与仅隐式化线性刚性部分（策略 B）两种分裂，推导了 $2\times2$ 放大矩阵，并通过谱半径扫描得到：策略 A 在 $\CFL\le1$ 时稳定；策略 B 的容许 CFL 上限为 $1/3$，且该限制不随 $\varepsilon\to0$ 而消失。数值实验验证了策略 B 在 $\CFL=0.7$ 时的失稳。

本文的方法论价值在于：把“容许分裂”“双变量稳定性函数”“耦合刚性极限”“L-稳定”这套语言，落实为一个可以计算、可以作图、可以用数值实验检验的完整流程。代码采用 Python 实现，结构清晰，可直接迁移到其他双曲守恒律与松弛系统。

\section{论文结构}

全文共七章，安排如下。

第一章为绪论，介绍研究背景、国内外现状、本文工作与结构。

第二章为预备知识，回顾常微分方程初值问题的基本算法与刚性概念，介绍 A-稳定、L-稳定、von Neumann 分析、有限体积方法与直线法，并给出 IMEX-Euler 的推导与双变量稳定性函数的初步定义。

第三章建立模型与分裂族，推导粘性 Burgers 方程的有限体积半离散系统，引入参数化分裂族，说明 IMEX-Euler 的实现细节与离散参数。

第四章是稳定性分析的核心，定义双变量稳定性函数与容许分裂，证明 $\theta\le1/2$ 的充要条件，讨论 L-稳定与离散符号修正。

第五章为数值验证，介绍实验设计，给出刚性相图、最大稳定步长、激波精度与收敛性结果，并讨论稳定性与精度的关系。

第六章为扩展，研究 Jin--Xin 松弛系统的两种源项分裂，推导放大矩阵，比较刚性极限行为，并做数值对比。

第七章总结全文，指出局限性与后续研究方向。

附录给出核心代码片段与复现说明。
